\section{Messaggi chiave}

\subsection{Milestone 1}

\begin{itemize}
    \item La costruzione di un dataset per la predizione della bugginess richiede numerose
    decisioni soggettive; documentarle esplicitamente è fondamentale per la riproducibilità.

    \item L'utilizzo delle FV come AV in assenza di AV esplicite riduce il ricorso a
    Proportion, preservando in parte l'informazione originale del ticket, tenendo sempre
    conto che si tratti di un'ipotesi verificabile solo parzialmente.

    \item I code smell sono shiftati di una release per riflettere la loro natura causale:
    i smell introdotti nella release $N$ non impattano quella stessa release, ma le successive.

    \item È importante non avvicinarsi temporalmente troppo alle release più recenti, in
    quanto possibilmente affette da bug dormienti (\textit{snoring bugs}).
\end{itemize}

\subsection{Milestone 2}

\begin{itemize}
    \item RandomForest domina su AUC e NPofB20 indipendentemente dal balancing: le
    probabilità calibrate per averaging su 100 alberi lo rendono il classificatore più
    affidabile per la prioritizzazione delle code review.

    \item La FS con soglia $0{,}0$ è un no-op su OpenJPA (tutte le 19 feature
    informative). I wrapper migliorano NaiveBayes attenuando la violazione dell'indipendenza,
    ma sono neutri per RF che gestisce la ridondanza internamente.

    \item Il bilanciamento introduce un trade-off irriducibile: Undersampling massimizza
    il Recall (meno bug sfuggiti, più falsi allarmi); SMOTE offre il miglior equilibrio
    su RF (Kappa $0{,}61$, NPofB20 $0{,}88$).

    \item NaiveBayes è strutturalmente limitato dalla correlazione tra LOC\_Touched, Churn
    e LOC\_Added --- il bottleneck non è il balancing, ma l'assunzione di indipendenza
    condizionale violata.

    \item La 10$\times$10 CV rinuncia alla correttezza temporale (temporal leakage) in
    favore di una stima stabile e confrontabile con poche release disponibili.
\end{itemize}

\subsection{Milestone 3}

\begin{itemize}
    \item L'impatto dei code smell è limitato: azzerare NSmells previene solo il 5,60\%
    dei bug predetti in B+, contro il 42\% del paper (Falessi). La differenza dipende
    dallo strumento (PMD vs SonarQube), dal contesto (open source vs industriale) e dallo
    sfasamento temporale della metrica.

    \item NSmells è ridondante rispetto alle metriche di processo ($\rho = 0{,}78$ con
    LOC, $> 0{,}65$ con Churn, LOC\_Touched): il ruleset \texttt{codestyle} introduce
    una correlazione artificiale con la dimensione del file.

    \item Lo studio di ablazione conferma la ridondanza: rimuovere NSmells lascia Kappa
    e AUC invariati ($\Delta \approx 0$), con un lieve miglioramento di Precision
    ($+0{,}004$). La feature è rumorosa, non informativa in modo indipendente.

    \item Il principale driver di bugginess in Apache OpenJPA è la storia di modifica
    del codice (LOC\_Touched, NR, Churn, Nfix), non la qualità strutturale misurata
    tramite PMD. Metriche di smell architetturali (God Class, Feature Envy, rilevabili
    con SonarQube) potrebbero fornire un segnale più indipendente.
\end{itemize}

\subsection{Milestone 4}
